% © 2026 Ing. David Jaroš — CC BY-NC-ND 4.0
%
% This work is licensed under a Creative Commons Attribution-NonCommercial-NoDerivatives
% 4.0 International License (CC BY-NC-ND 4.0).
%
% License History: Earlier drafts (up to v0.3) were released under CC BY 4.0.
% From v0.4 onward, all material is released under CC BY-NC-ND 4.0 to protect
% the integrity of the theoretical work during ongoing academic development.

% UBT-AI-PROVENANCE-BEGIN
% schema: ubt-ai-provenance/v1
% tier: C_working
% ai_assistance: disclosed
% human_review: risk-based
% editorial_responsibility: Ing. David Jaroš
% policy: ../../AI_PROVENANCE.md
% notice: Working material; exhaustive human review is not claimed.
% UBT-AI-PROVENANCE-END

\documentclass[11pt,a4paper]{article}
\usepackage{amsmath,amssymb,amsthm}
\usepackage{geometry}
\usepackage{hyperref}
\geometry{margin=2.5cm}

\newtheorem{proposition}{Proposition}[section]
\newtheorem{remark}[proposition]{Remark}

\title{QED Sector from \texorpdfstring{$S[\Theta]$}{S[Theta]} in UBT}
\author{Ing.\ David Jaroš}
\date{2026-05-10}

\begin{document}
\maketitle

\section{Starting point: canonical UBT action}
The canonical action is
\begin{equation}
S[\Theta]=\int d^4x\,\sqrt{-g}\left[
\mathrm{Tr}\big((D_\mu\Theta)^\dagger(D^\mu\Theta)\big)
-\frac14\,W^i_{\mu\nu}W^{i\mu\nu}
-\frac14\,B_{\mu\nu}B^{\mu\nu}
+\cdots\right],
\end{equation}
with
\begin{equation}
D_\mu=\partial_\mu+ig\,\tau^iW^i_\mu+ig'YB_\mu+\cdots.
\end{equation}
This is the unique canonical entry point for the electromagnetic sector.

\section{Identification of the U(1) sector from \texorpdfstring{$\Theta$}{Theta}}
In UBT, the right phase action
\begin{equation}
\Theta(x)\mapsto \Theta(x)e^{i\chi(x)}
\end{equation}
defines the abelian connection sector. At electroweak projection,
\begin{equation}
A_\mu=\sin\theta_W\,W^3_\mu+\cos\theta_W\,B_\mu,
\qquad
e=g\sin\theta_W=g'\cos\theta_W,
\end{equation}
and the unbroken generator is $Q=T_3+Y/2$.
Therefore the low-energy covariant derivative is
\begin{equation}
D^{\mathrm{EM}}_\mu=\partial_\mu+i e\,Q\,A_\mu.
\end{equation}

\section{Low-energy QED-like action}
After integrating out heavy modes and keeping the photon plus light charged modes,
\begin{equation}
S_{\mathrm{eff}}=\int d^4x\,\sqrt{-g}\left[
\sum_r \big|\big(\nabla_\mu+i e q_r A_\mu\big)\psi_r\big|^2
-\frac14 F_{\mu\nu}F^{\mu\nu}
+\cdots\right],
\end{equation}
with $F_{\mu\nu}=\partial_\mu A_\nu-\partial_\nu A_\mu$.
Thus UBT does generate the QED gauge kinetic structure in the infrared.

\section{Isolation of the \texorpdfstring{$(1/4e^2)F^2$}{(1/4e^2)F^2} form}
Define a charge-normalized connection
\begin{equation}
\mathcal{A}_\mu:=eA_\mu,
\qquad
\mathcal{F}_{\mu\nu}=\partial_\mu\mathcal{A}_\nu-\partial_\nu\mathcal{A}_\mu=eF_{\mu\nu}.
\end{equation}
Then
\begin{equation}
D^{\mathrm{EM}}_\mu=\partial_\mu+i q_r\mathcal{A}_\mu,
\qquad
-\frac14 F_{\mu\nu}F^{\mu\nu}=-\frac{1}{4e^2}\mathcal{F}_{\mu\nu}\mathcal{F}^{\mu\nu}.
\end{equation}
Hence the coefficient required by the task is present as a normalization-equivalent form
of the canonical gauge kinetic term.

\begin{remark}[Current status]
What is fixed by current UBT documents is the existence of the abelian gauge sector and
its canonical kinetic form. What is not yet fixed from first principles is the absolute
normalization value of $e$ itself.
\end{remark}

\end{document}
